% title: 课程说明
% nav_title:00-课程说明
% author：潘宇琦
% summary: 
\documentclass[12pt, a4paper, oneside]{ctexart}
\usepackage{amsmath,amsthm,amssymb,graphicx}
\usepackage{float}
\usepackage[margin=2.5cm]{geometry}
\usepackage{xcolor}
\usepackage{enumitem}
\usepackage{hyperref}
\hypersetup{
    colorlinks=true,
    linkcolor=blue,
    urlcolor=blue,
    citecolor=blue,
}


\begin{document}
\section*{个人看法}

线性代数是大家上大学接触到的第一门比较有难度的数学课。话虽然这样讲，但是拿分还是很容易的，你只需要买一本由张天德编著的《线性代数习题精选精解》，反复地刷里面的题目，基本上就能保证你最后获得很高的分数。刷分可以这样刷，但是学习不能这么学对吧...

但是我想说明的是，我们的课本《线性代数(第三版)》是一本足够清楚明白的课本，我觉得我就算更新文档也更新不出什么新的地方，所以线性代数的文档更新暂时被挂起。

\section*{关于文档的参考}

\textbf{本文档的所有知识点几乎全部来自于山东大学数学学院的《线性代数(第三版)》，有部分习题也来自于这本书}，这也是24级这门课程的课本。

除了这本课本之外，我们还参考了下面几本教材：
\begin{itemize}
    \item 李尚志《线性代数》
    \item 杨子胥《高等代数》
\end{itemize}

\section*{关于文档所涉及的章节}

这本课本只有五章的内容，其实每一章都很重要，所以文档会详细解释每一章的内容。

\section*{其他}

文档的PDF版以及其他电子资料可以前往仓库下载。



\end{document}